%============================================================
% Chapter 27: Games on Arrays
% STATUS: PLANNED  (generated by tools/outline.py; replace this file when the chapter is written)
%============================================================
\chapter{Games on Arrays}
\label{ch:games}

\begin{tcolorbox}[enhanced, colback=white, colframe=secondaryblue, boxrule=1pt, arc=2mm]
\textbf{\large\textcolor{primaryblue}{Two players, one array: find the winning invariant --- parity, symmetry, the first decisive element, or a Nim sum --- instead of searching the game tree.}}
\end{tcolorbox}

\chapterfacts{1100}{2300}{5}{10}{30}

\begin{plannedbox}[Planned Contents --- this chapter is not written yet]
\textbf{Part A --- Lessons}
\begin{enumerate}[leftmargin=*, label=\textbf{A\arabic*.}, itemsep=1pt]
    \item Winning and losing positions; why the game tree is too large to search
    \item Parity and counting arguments
    \item Symmetry and mirror strategies
    \item Games decided by the first element that gives a real choice
    \item Nim and the XOR rule as a tool (Sprague--Grundy theory is the math notes; interval games are the DP book)
\end{enumerate}
\textbf{Part B --- Patterns} (with candidate problems, checked when the chapter is written)
\begin{enumerate}[leftmargin=*, label=\textbf{P\arabic*.}, itemsep=1pt]
    \item \textbf{Parity arguments}
    \item \textbf{Mirror strategies}
    \item \textbf{Decided by the first real choice} --- Sequential Nim {\small\color{darkgray}(CF 1382B)}
    \item \textbf{Nim sums on arrays}
    \item \textbf{Games on operations}
\end{enumerate}
\end{plannedbox}
